%!TEX root = 1_a_Generalised IIT.tex
\section{Systems} \label{sec:sys}

The first step in defining an Integrated Information Theory (IIT) is to specify a class $\Sys$ of physical \emph{systems} to be studied. Each element $S \in \Sys$ is interpreted as a model of one particular physical system. In order to apply the IIT algorithm, it is only necessary that each element $S$ come with the following features.

\begin{definition} \label{def:sys-class}
A \emph{system class} $\Sys$ is a class each of whose elements $S$, called \emph{systems}, come with the following data:
\begin{enumerate}[label=\arabic*.]
\item a set $\St(S)$ of \emph{states};
\item for every $s \in \St(S)$, a set $\Sub_s(S) \subset \Sys$ of subsystems and for each $M \in \Sub_s(S)$ an
induced state $s|_M \in \St(M)$;
\item \label{it:Decomp} a set $\mathbb D_S$ of {\em decompositions}, with a given \emph{trivial decomposition} $1 \in \mathbb D_S$;
\item for each $z \in \mathbb D_S$ a corresponding \emph{cut system} $\cut{S}{z} \in \Sys$ and for each state $s \in \St(S)$ a corresponding \emph{cut state} $\cut{s}{z} \in \St(\cut{S}{z})$.
\end{enumerate}
\end{definition}
Moreover, we require that $\Sys$ contains a distinguished \emph{empty system}, denoted~$I$, and that $I \in \Sub(S)$ for all $S$. For the IIT algorithm to work, we need to assume furthermore that the number of subsystems remains the same under cuts or changes of states,
i.e. $\Sub_s(S) \simeq \Sub_{s'}(S)$ for all $s, s' \in \St(S)$ and $\Sub_s(S) \simeq \Sub_{\cut{s}{z}}(\cut{S}{z})$ for all $z \in \mathbb D_S$. Here, $\simeq$ indicates bijections. Note that subsystems of a system may depend on the state of the system, in accordance with classical IIT. 

In this article we will assume that each set $\Sub_s(S)$ is finite, discussing the extension to the infinite case in Section \ref{sec:discussion}. We will give examples of system classes and for all following definitions in Sections~\ref{sec:Classical} and~\ref{sec:quantum-IIT}.



\section{Experience} \label{sec:Experience}

An IIT aims to specify for each system in a particular state its \emph{conscious experience}. As such, it will require a mathematical model of such experiences. Examining classical IIT, we find the following basic features of the final experiential states it describes which are needed for its algorithm.

Firstly, each experience $e$ should crucially come with an \emph{intensity}, given by a number $\norm e \norm$ in the non-negative reals  $\R^+$ (including zero). 
This intensity will finally correspond to the overall intensity of experience, usually denoted by $\Phi$. Next, in order to compare experiences, we require a notion of \emph{distance} $d(e,e')$ between any pair of experiences $e, e'$. Finally, the algorithm will require us to be able to \emph{rescale} any given experience $e$ to have any given intensity. Mathematically, this is most easily encoded by letting us multiply any experience $e$ by any number $r \in \R^+$.
In summary, a minimal model of experience in a generalised IIT is the following.

\begin{definition} \label{def:exp-space}
An \emph{experience space} is a set $E$ with: 
\begin{enumerate}[label=\arabic*.]
\item 
an \emph{intensity} function $\norm . \norm \colon E \to \R^+$;
\item 
a \emph{distance} function $d \colon E \times E \to \mathbb{R}^+$;
\item 
a \emph{scalar multiplication} $\R^+ \times E \to E$, denoted $(r,e) \mapsto r \cdot e$, satisfying
\begin{align*}
\norm r \cdot e \norm = r \cdot \norm e \norm \qquad 
r \cdot (s \cdot e) = (rs) \cdot e \qquad
1 \cdot e = e
\end{align*}
for all $e \in E$ and $r, s \in \R^+$.
\end{enumerate}
\end{definition}


We remark that this same axiomatisation will apply both to the full space of experiences of a system, as well as to the the spaces describing components of the experiences (`concepts' and `proto-experiences' defined in later sections). We note that the distance function does not necessarily have to satisfy the axioms of a metric. While this and further natural axioms such as $d(r \cdot e, r \cdot f) = r \cdot d(e,f)$ might hold, they are not necessary for the IIT algorithm.

The above definition is very general, and in any specific theory the experiences may come with richer further structure. The following example describes the experience space used in classical IIT.

\begin{example} \label{ex:IITProto}
Any metric space $(X,d)$ may be extended to an experience space $\bar X := X \times \mathbb{R}^+$ in various ways. E.g., 
one can define $\norm (x,r) \norm = r$, $r \cdot (x,s) = (x, rs)$ and define the distance as
\begin{equation}\label{eq:ExIITProto}
d\big((x,r),(y,s)\big) = r \, d(x,y) 
\end{equation}
This is the definition used in classical IIT (cf. Section~\ref{sec:Classical} and Appendix~\ref{sec:comparison}).
\end{example}


An important operation on experience spaces is taking their \emph{product}. 

\begin{definition} \label{def:prod-spaces}
For experience spaces $E$ and $F$, we define the product to be the space $E \times F$ with distance
\begin{equation}\label{eq:DistanceProduct}
d\big( (e,f) , (e',f') \big) = d(e,e') + d(f, f')  \:,
\end{equation} 
intensity $\norm (e,f) \norm = \max \{ \norm e \norm, \norm f \norm \}$ and scalar multiplication $r \cdot (e,f) = (r \cdot e, r \cdot f)$.
This generalises to any finite product $\prod_{i \in I} E_i$ of experience spaces.
\end{definition}





\section{Repertoires} \label{sec:repertoires}

% The main aim of an IIT is to assign an experience to each system~$S$ in a state $s \in \St(S)$, each of which is an element of an experience space. 

In order to define the experience space and individual experiences of a system $S$, an IIT utilizes basic building blocks called `repertoires', which we will now define.
Next to the specification of a system class, this is the essential data necessary for the IIT algorithm to be applied.

Each repertoire describes a way of `decomposing' experiences, in the following sense. Let $D$ denote any set with a distinguished element $1$, for example the set $\mathbb D_S$ of decompositions of a system $S$, where the distinguished element is the trivial decomposition $1 \in \mathbb D_S$.

\begin{definition}\label{def:Decomp}
Let $e$ be an element of an experience space $E$. A \emph{decomposition of $e$ over $D$} is a mapping $\bar{e} \colon D \to E$ with $\bar{e}(1) = e$.
\end{definition}

In more detail, a repertoire specifies a proto-experience for every pair of subsystems and describes how this experience changes if the subsystems are decomposed.
This allows one to assess how integrated the system is with respect to a particular repertoire. Two repertoires are necessary for the IIT algorithm to be applied, together called the cause-effect repertoire.

For subsystems $M,P \in \Sub_s(S)$, define $\mathbb D_{M,P} := \mathbb D_M \times \mathbb D_P$. This set describes the decomposition of both subsystems simultaneously. It has a distinguished element $1 = (1_M, 1_P)$.

\begin{definition} \label{def:repertoire}
A \emph{cause-effect repertoire} at $S$ is given by a choice of experience space $\PExp(S)$, called the space of \emph{proto-experiences}, and for each 
$s \in \St(S)$ and $M, P \in \Sub_s(S)$, a pair of elements 
\begin{equation}\label{eq:CauseEffectRep}
\caus_s(M,P) \  ,\  \eff_s(M,P) \  \in \ \PExp(S)
\end{equation}
and for each of them a decomposition over $\mathbb D_{M,P}$.
\end{definition}

Examples of cause-effect repertoires will be given in Sections~\ref{sec:Classical} and~\ref{sec:quantum-IIT}. A general definition in terms of process theories is given in~\cite{IITProcess}. For the IIT algorithm, a cause-effect repertoire needs to be specified for every system $S$, as in the following definition.

\begin{definition}\label{def:cestructure}
A \emph{cause-effect structure} is a specification of a cause-effect repertoire for every $S \in \Sys$ such that
\begin{equation}\label{eq:cestructure}
\PExp(S) = \PExp(\cut{S}{z}) \quad \textrm{ for all } \quad z \in \mathbb D_S \:.
\end{equation}
%together with an embedding $\PExp(M) \hookrightarrow \PExp(S)$, for all subsystems $M$ of $S$.
\end{definition}

The names `cause' and `effect' highlight that the definitions of $\caus_s(M,P)$ and  $\eff_s(M,P)$ in classical and quantum IIT describe the causal dynamics of the system. More precisely, they are intended to capture the manner in which the `current' state~$s$ of the system, when restricted to $M$, constrains the `previous' or `next' state of $P$, respectively. 


\section{Integration} \label{sec:Integration}

We have now introduced all of the data required to define an IIT; namely, a system class along with a cause-effect structure. From this, we will give an algorithm aiming to specify the conscious experience of a system. Before proceeding to do so, we introduce a conceptual short-cut
which allows the algorithm to be stated in a concise form. This captures the core ingredient of an IIT, namely the computation of how integrated an entity is.

\begin{definition}\label{def:IntScaling}
Let $E$ be an experience space and $e$ an element with a decomposition over some set $D$. The \emph{integration level} of $e$ relative to this decomposition is 
\begin{equation}\label{eq:IntegrLevel}
\phi(e) := \min_{1 \neq z \in D} d(e, \bar{e}(z)) \:.
\end{equation}
Here, $d$ denotes the distance function of $E$, and the minimum is taken over all elements of $D$ besides $1$. The \emph{integration scaling} of $e$ is then the element of $E$ defined by
\begin{equation}\label{eq:IntegrData}
	\intscaling(e) := \phi(e) \cdot \hat{e} \:,
\end{equation}
where $\hat e$ denotes the \emph{normalization} of $e$, defined as
\begin{align*}
\hat e := \begin{cases}
\frac{1}{\norm e \norm} \cdot e & \textrm{ if } \norm e \norm \neq 0 \\
 e & \textrm{ if } \norm e \norm = 0 \:.
\end{cases}
\end{align*}
Finally, the \emph{integration scaling} of a pair $e_1, e_2$ of such elements is the pair
\begin{equation}\label{eq:IntScalingPair}
\intscaling(e_1, e_2) :=
(\phi \cdot \hat{e_1},
\phi \cdot \hat{e_2})
\end{equation}
where $\phi := \min(\phi(e_1), \phi(e_2))$
is the minimum of their integration levels.
\end{definition}

We will also need to consider indexed collections of decomposable elements.
Let $S$ be a system in a state $s \in \St(S)$ and assume that for every $M \in \Sub_s(S)$
an element $e_M$ of some experience space $E_M$ with a decomposition over some $D_M$ is given. We call
$(e_M)_{M \in \Sub_s(S)}$ a \emph{collection} of decomposable elements, and denote it as
$(e_M)_M$.

\begin{definition}\label{def:IntScalingColl}
The \emph{core} of the collection $(e_M)_{M}$ is the subsystem $C \in \Sub(S)$ for which $\phi(e_C)$ is maximal.%
\footnote{If the maximum does not exist, we define the core to be the empty system $I$.}
The \emph{core integration scaling} of the collection is $\intscaling(e_C)$.
The \emph{core integration scaling} of a pair of collections $(e_M, f_M)_{M}$ is
$\intscaling(e_{C}, f_{D})$, where $C, D$ are the cores of $(e_M)_M$ and $(f_M)_M$, respectively.
\end{definition}


\section{Constructions - Mechanism Level} \label{sec:alg-mech}

Let $S \in \Sys$ be a physical system whose experience in a state $s \in \St(S)$ is to be determined.
The first level of the algorithm involves fixing some subsystem $M \in \Sub_s(S)$, referred to as a `mechanism', and associating to it an object called its `concept' which belongs to the \emph{concept space}
\begin{equation}\label{eq:Concepts}
\Concepts(S) := \PExp(S) \times \PExp(S) \:.
\end{equation}


For every choice of $P \in \Sub_s(S)$, called a `purview', the repertoire values $\caus_s(M,P)$ and $\eff_s(M,P)$ are elements of $\PExp(S)$ with given decompositions over $\mathbb{D}_{M,P}$. Fixing $M$, they form collection of decomposable elements,
\begin{align}\begin{split}\label{eq:CER}
\caus_s(M) &:= (\caus_s(M,P))_{P \in \Sub(S)} \\
\eff_s(M) &:= (\eff_s(M,P))_{P \in \Sub(S)} \:.
\end{split}\end{align}


The \emph{concept} of $M$ is then defined as the core integration scaling of this pair of collections,
% each assignment,
\begin{equation}\begin{split}\label{eq:DefConcept}
\con{S,s}{M} := \textrm{Core integration scaling of }  \left(\caus_s(M), \eff_s(M) \right) \:.
\end{split}\end{equation}
It is an element of $\Concepts(S)$. Unravelling our definitions, the concept thus consists of the values of the cause and effect repertoires at their respective `core' purviews $P^c, P^e$, i.e.~those which make them `most integrated'. These values $\caus(M,P^c)$ and $\eff(M,P^e)$ are then each rescaled to have intensity given by the minima of their two integration levels.

\section{Constructions - System Level} \label{sec:alg-sys}

The second level of the algorithm specifies the experience of the system $S$ in state $s$.
To this end, all concepts of a system are collected to form its \emph{Q-shape}, defined as
\begin{equation} \label{eq:Q-Shape}
\QShape_s(S) := (\con{S,s}{M})_{M \in \Sub_s(S)}  \:.
\end{equation}
This is an element of the space 
\begin{equation}\label{eq:ExpS}
\Exp(S) = \Concepts(S)^{n(S)} \:,
\end{equation}
where $n(S) := |\Sub_s(S)|$, which is finite and independent of the state $s$ according to our assumptions.
We can also define a Q-shape for any cut of $S$. Let $z \in \mathbb D_S$ be a decomposition, $\cut{S}{z}$ the corresponding cut system and $\cut{s}{z}$ be the
corresponding cut state. We define
\begin{equation} \label{eq:Q-Shape-cuts}
\QShape_s(\cut{S}{z}) := (\con{\cut{S}{z},\cut{s}{z}}{M})_{M \in \Sub_{\cut{s}{z}}(\cut{S}{z})} \:.
\end{equation}
Because of~\eqref{eq:cestructure}, and since the number of subsystems remains the same when cutting, 
$\QShape_s(\cut{S}{z})$ is also an element of $\Exp(S)$.
This gives a map
\begin{align*}
\bar \QShape_{S,s}: \mathbb D_S &\rightarrow \Exp(S)\\
z &\mapsto \QShape_s(\cut{S}{z})
\end{align*}
which is a decomposition of $\QShape_s(S)$ over $\mathbb D_S$. Considering this map for every subsystem of $S$ gives a collection of decompositions
defined as
\begin{align*}
\QShape(S,s) := \big( \bar \QShape_{M,s|_M} \big)_{M \in \Sub_s(S)} \:
\end{align*}
This is the system level-object of relevance and is what specifies the experience of a system according to IIT.

\begin{definition}\label{def:ActualExp}
The actual \emph{experience} of the system $S$ in the state $s \in \St(S)$ is
\begin{equation}\label{eq:Exp}
\Exp(S,s) :=  \textrm{Core integration scaling of } \QShape(S,s) \:.
\end{equation}
\end{definition}
The definition implies that $\Exp(S,s) \in \Exp(M)$, where $M \in \Sub_s(S)$ is the core of the collection $\QShape(S,s)$, called the \emph{major complex}.
It describes which part of the system $S$ is actually conscious.
%Due to the last condition in Definition~\ref{def:cestructure}, we have $\Concepts(M)$
In most cases there will be a natural embedding $\Exp(M) \hookrightarrow \Exp(S)$ for a subsystem $M$ of $S$, allowing us to view $\Exp(S,s)$ as an element of $\Exp(S)$ itself. Assuming this embedding
to exist allows us to define an Integrated Information Theory concisely in the next section.

\section{Integrated Information Theories}\label{sec:IITs}

We can now summarise all that we have said about IITs.

\begin{definition}
An \emph{Integrated Information Theory} is determined as follows. The \emph{data} of the theory is a system class $\Sys$ along with a cause-effect structure. The theory then gives a mapping
% An \emph{integrated information theory} is a mapping 
\begin{equation} \label{eq:IIT-High-Level}
\begin{tikzcd}
\Sys\rar{\Exp} & \Expcat 
\end{tikzcd}
\end{equation}
into the class $\Expcat$ of all experience spaces, sending each system $S$ to its space of experiences $\Exp(S)$ defined in~\eqref{eq:ExpS}, and a mapping
\begin{align}\begin{split} \label{eq:IIT-On-States}
\St(S) &\to \Exp(S) \\ 
s &\mapsto \Exp(S,s) 
\end{split}\end{align}
which determines the experience of the system when in a state $s$, defined in~\eqref{eq:Exp}.

\noindent The \emph{quantity} of the system's experience is given by 
\begin{equation*}
\Phi(S,s) := \norm \Exp(S,s) \norm \:,
\end{equation*}
and the quality of the system's experience is given by the normalized experience $\hat \Exp(S,s)$.
The experience is located in the core of the collection $\QShape(S,s)$, called \emph{major complex}, which is a subsystem of $S$.
\end{definition}


In the next sections we specify the data of several example IITs.


